\documentclass[10pt,twocolumn,a4paper]{article}
\usepackage[T1]{fontenc}
\usepackage[utf8]{inputenc}
\usepackage{lmodern}
\usepackage{booktabs}
\usepackage[hidelinks]{hyperref}
\title{Heurísticas de layout na conversão de PDF}
\author{Equipe markDOWN}
\date{Novembro de 2026}
\begin{document}
\maketitle

\section{Introdução}
Documentos produzidos com \LaTeX{} são comuns na academia. Eles usam fontes próprias, hifenização automática e, com frequência, duas colunas. Essas características tornam a extração de texto um desafio interessante para qualquer conversor que pretenda reconstruir a estrutura original do documento com fidelidade.

A hifenização, em particular, quebra palavras no fim da linha. Um bom conversor precisa juntar de volta essas partes para que a busca e a leitura funcionem no resultado final.

\section{Abordagem}
O conversor trabalha em etapas independentes:
\begin{enumerate}
  \item agrupa os trechos de texto em linhas;
  \item separa as colunas da página;
  \item identifica títulos, listas e tabelas;
  \item gera o Markdown.
\end{enumerate}

Cada etapa é uma função pura, o que facilita os testes. O código está disponível em \href{https://github.com/joaopedrogurgel/markDOWN}{github.com/joaopedrogurgel/markDOWN}.

\subsection{Estilos}
Trechos em \textbf{negrito} e em \emph{itálico} são reconhecidos pelo nome da fonte embutida no arquivo.

\section{Resultados}
A tabela abaixo resume os tipos de documento avaliados.

\begin{tabular}{lrr}
\toprule
Tipo & Arquivos & Páginas \\
\midrule
Artigos & 12 & 96 \\
Relatórios & 8 & 40 \\
Memorandos & 10 & 14 \\
\bottomrule
\end{tabular}

\section{Conclusão}
Heurísticas simples, bem testadas, resolvem a maior parte dos documentos de texto. Os casos difíceis, como fórmulas e tabelas complexas, ficam para trabalhos futuros.

\end{document}
